\documentclass[preprint]{revtex4-2}
\usepackage{amsmath}
\usepackage{graphicx}
\usepackage{comment}
\newcommand{\aap}{Astron. Ap.}
\newcommand{\apjs}{Astrophys. J. Suppl.}
\newcommand{\gca}{Geochim. Cosmochim. Acta}
\newcommand{\epsl}{Earth Planet. Sci. Lett.}

\begin{document}

\title{Mixing Relations\footnote{Supplemental material for the paper
{\it Constraints on the building blocks of the Solar System: Insights from nucleosynthetic isotope anomalies}, XXX (2023), authors: K. Bermingham,
B. S. Meyer, and K. Mezger, submitted}}
\author{Bradley S. Meyer}
\affiliation{Department of Physics and Astronomy, Clemson University, Clemson,
SC 29634, USA}
\author{Katherine R. Bermingham}
\affiliation{ Department of Earth and Planetary Sciences, Rutgers University, Piscataway, NJ 08854, USA}
\author{Klaus Mezger}
\affiliation{Institut f\"ur Geologie, Baltzerstrasse 1+3 3012 Bern, Schweiz}

\begin{abstract}
Mixing relations for the NRLEE paper.
\end{abstract}

\maketitle

\section{Introduction}\label{sec:intro}

Derive the mixing relations.

\section{Relations for Different Standards} \label{sec:initial}

It is convenient to consider anomalies measured  against different standards.
Suppose $X_i(s)$ is the mass fraction of species $i$ in sample $s$, that is,
$X_i(s)$ is the number of grams of species $i$ in sample $s$ relative to the
total number of grams in $s$.  Let
\begin{equation}
\mu_{ij}^{(a)}(s) = 10^6 \left(\frac{X_i(s)/X_i(a)}{X_j(s)/X_j(a)} - 1\right)
\label{eq:mu_ija}
\end{equation}
be the anomaly in species $i$ relative to species $j$ in sample $s$
measured against standard $a$ in parts per million.
If $b$ is a second possible standard, then
\begin{equation}
\mu_{ij}^{(b)}(s) - \mu_{ij}^{(a)}(s) =
10^6 \left(\frac{X_i(s)/X_i(b)}{X_j(s)/X_j(b)} -
\frac{X_i(s)/X_i(a)}{X_j(s)/X_j(a)}
\right)
\label{eq:diff_mu}
\end{equation}
This may be written
\begin{equation}
\mu_{ij}^{(b)}(s) - \mu_{ij}^{(a)}(s) =
10^6 \frac{X_i(s)}{X_j(s)}\frac{X_j(a)}{X_i(a)}
\left(\frac{X_i(a)/X_i(b)}{X_j(a)/X_j(b)} - 1\right)
\label{eq:diff_mu2}
\end{equation}
From the definition of $\mu_{ij}$, this becomes
\begin{equation}
\mu_{ij}^{(b)}(s) - \mu_{ij}^{(a)}(s) =
\mu_{ij}^{(b)}(a) \left(1 + \frac{\mu_{ij}^{(a)}(s)}{10^6}\right)
\label{eq:diff_mu3}
\end{equation}
which may be rearranged to
\begin{equation}
\mu_{ij}^{(b)}(s) =
\mu_{ij}^{(a)}(s) \left(1 + \frac{\mu_{ij}^{(b)}(a)}{10^6}\right)
+ \mu_{ij}^{(b)}(a)
\label{eq:diff_mu4}
\end{equation}
With this last equation, an anomaly relative to standard $b$ can be
related to the anomaly relative to standard $a$ given knowledge of the
abundance offset $\mu_{ij}^{(b)}(a)$ between the two standards.
The corresponding relation for anomalies measured in per mil is
\begin{equation}
\delta_{ij}^{(b)}(s) =
\delta_{ij}^{(a)}(s) \left(1 + \frac{\mu_{ij}^{(b)}(a)}{10^3}\right)
+ \mu_{ij}^{(b)}(a)
\label{eq:diff_delta}
\end{equation}

\section{Reservoir Mixing}

Suppose a sample $s$ results from the mixture of materials from two reservoirs
1 and 2.  The resulting mass of species $i$ in $s$ can be written
\begin{equation}
M_i(s) = g_i(1, s) X_i(1) M(1, s) + g_i(2, s) X_i(2) M(2, s)
\label{eq:mix1}
\end{equation}
$M(1, s)$ and $M(2, s)$ are the masses of reservoirs 1 and 2 contributing to the
mixture, respectively, $X_i(r)$, with $r = 1,2$ are the mass fractions
of species $i$ in reservoirs 1 and 2, and $g_i(1, s)$ and $g_i(2, s)$ are
numbers
between 0 and 1 allowing for incomplete addition of reservoir materials to
the budget of species $i$ in sample $s$.  The total mass of sample $s$ is
the sum of the masses of all species; thus,
\begin{equation}
M(s) = G(1, s) M(1, s) + G(2, s) M(2, s)
\label{eq:mix2}
\end{equation}
where
\begin{equation}
G(r, s) = \sum_i g_i(r, s) X_i(r)
\label{eq:G}
\end{equation}
for $r = 1, 2$.  Since the sum of all mass fractions in a reservoir is
unity, $G(r, s) = 1$ if $g_i(r, s) = 1$ for all $i$.

It is convenient to define $M(s) = M(1, s) + M(2, s)$ and
$M(1, s) = f_s M(s)$.  Since
$X_i(s) = M_i(s) / M(s)$, this results in the equation
\begin{equation}
X_i(s) = \frac{f_s g_i(1, s) X_i(1) + \left(1 - f_s\right) g_i(2, s) X_i(2)}{f_s G(1, s) +
\left(1 - f_s\right) G(2, s)}
\label{eq:XG}
\end{equation}
Again, if $g_i(r, s) = 1$ for all $i$, this reduces to the well-known result
\begin{equation}
X_i(s) = f_s X_i(1) + \left(1 - f_s\right) X_i(2)
\label{eq:Xmix}
\end{equation}
From Eq. (\ref{eq:mu_ija}), an anomaly in sample $s$
in species $i$ relative to species
$j$ referenced to standard $a$ is given by
\begin{equation}
\mu_{ij}^{(a)}(s) =
10^6 \left(\frac{f_s \left[g_i(1, s) \frac{X_i(1)}{X_i(a)} -
g_j(1, s) \frac{X_j(1)}{X_j(a)}\right]
+
\left(1 - f_s\right) \left[g_i(2, s)\frac{X_i(2)}{X_i(a)}
- g_j(2, s)\frac{X_j(2)}{X_j(a)}\right]}
{f_s g_j(1, s) \frac{X_j(1)}{X_j(a)} +
\left(1 - f_s\right) g_j(2,s)\frac{X_j(2)}{X_j(a)}}\right)
\label{eq:mu_g}
\end{equation}
In the case that $a = 2$, Eq. (\ref{eq:mu_g}) can be written
\begin{equation}
\mu_{ij}^{(a)}(s) =
10^6 \left(\frac{f_s \left[g_i(1, s) \frac{X_i(1)}{X_i(a)} -
g_j(1, s) \frac{X_j(1)}{X_j(a)}\right]
+
\left(1 - f_s\right) \left[g_i(2, s) - g_j(2, s)\right]}
{f_s g_j(1, s) \frac{X_j(1)}{X_j(a)} + \left(1 - f_s\right) g_j(2,s)}\right)
\label{eq:mu_g2a}
\end{equation}
Furthermore,
if $g_i(r, s) = g_j(r, s)$, as would be expected if $i$ and $j$ are isotopes
of the same element, then Eq. (\ref{eq:mu_g2a}) becomes
\begin{equation}
\mu_{ij}^{(a)}(s) =
10^6 \left(\frac{f_s \left[\frac{X_i(1)}{X_i(a)} -
\frac{X_j(1)}{X_j(a)}\right]}
{f_s \frac{X_j(1)}{X_j(a)} + \left(1 - f_s\right) \frac{g_j(2,s)}{g_j(1, s)}}\right)
\label{eq:mu_g2a2}
\end{equation}
which reduces to the correct limits $\mu_{ij}^{(a)}(a) = 0$ for $f_s = 0$
and $\mu_{ij}^{(a)}(s) = \mu_{ij}^{(a)}(1)$ for $f_s = 1$.

\section{A Derived Standard}

If a known standard $a$ is the result of the mixture of two reservoirs 1 and 2,
it is possible to eliminate the mass fraction of species in reservoir 2
and express the resulting mixtures in terms of reservoir 1 and the standard.
By Eq. (\ref{eq:XG}), one may write
\begin{equation}
\frac{f_a G(1, a) + \left(1 - f_a\right) G(2, a)}{\left(1 - f_a\right) g_i(2, a)}
X_i(a)
= \frac{f_a g_i(1, a)}{\left(1 - f_a\right) g_i(2, a)} X_i(1) + X_i(2)
\label{eq:XGa}
\end{equation}
For sample $s$, one may similarly write
\begin{equation}
\frac{f_s G(1, s) + \left(1 - f_s\right) G(2, s)}{\left(1 - f_s\right) g_i(2, s)}
X_i(s)
= \frac{f_s g_i(1, s)}{\left(1 - f_s\right) g_i(2, s)} X_i(1) + X_i(2)
\label{eq:XGs}
\end{equation}
If one subtracts Eq. (\ref{eq:XGa}) from Eq. (\ref{eq:XGs}), one finds
\begin{multline}
\frac{f_s G(1, s) + \left(1 - f_s\right) G(2, s)}{\left(1 - f_s\right) g_i(2, s)}
X_i(s)
-
\frac{f_a G(1, a) + \left(1 - f_a\right) G(2, a)}{\left(1 - f_a\right) g_i(2, a)}
X_i(a)
= \\
X_i(1) \left(
\frac{f_s g_i(1, s)}{\left(1 - f_s\right) g_i(2, s)}
- \frac{f_a g_i(1, a)}{\left(1 - f_a\right) g_i(2, a)}\right)
\label{eq:XGdiff}
\end{multline}
This may be rewritten
\begin{multline}
\frac{X_i(s)}{X_i(a)}
=
\left[
\frac{(1 - f_s) g_i(2, s)}{(1 - f_a) g_i(2, a)}
+
A_i(a)\frac{X_i(1)}{X_i(a)}
\left\{
\frac{f_s g_i(1, s)}{f_a g_i(1, a)}
- \frac{\left(1 - f_s\right) g_i(2, s)}{\left(1 - f_a\right) g_i(2, a)}\right\}
\right]
\\
\left[\frac{f_a G(1, a) + (1 - f_a) G(2,a)}{f_s G(1, s) + (1 - f_s) G(2,s)}\right]
\label{eq:XGdiff2}
\end{multline}
with
\begin{equation}
A_i(a) = \frac{f_a g_i(1, a)}{f_a G(1, a) + (1 - f_a) G(2,a)}
\label{eq:A}
\end{equation}
It is convenient to write the mass fraction ratio
\begin{equation}
\frac{X_i(s)}{X_i(a)} = N(s, a) B_i(s, a)
\label{eq:Xi}
\end{equation}
where the normalization for sample $s$ and standard $a$ is
\begin{equation}
N(s, a) = \frac{f_a G(1, a) + (1 - f_a) G(2,a)}{f_s G(1, s) + (1 - f_s) G(2,s)}
\label{eq:norm}
\end{equation}
and the species specific quantity $B_i(s, a)$ is given by
\begin{equation}
B_i(s, a) = 
\frac{(1 - f_s) g_i(2, s)}{(1 - f_a) g_i(2, a)}
+
A_i(a) \frac{X_i(1)}{X_i(a)}
\left\{
\frac{f_s g_i(1, s)}{f_a g_i(1, a)}
- \frac{\left(1 - f_s\right) g_i(2, s)}{\left(1 - f_a\right) g_i(2, a)}\right\}
\label{eq:Bi}
\end{equation}
Analogous equations hold for species $j$.
Because $N(s, a)$ is species independent,
\begin{equation}
\mu_{ij}^{(a)} = 10^6 \left(\frac{B_i(s, a)}{B_j(s, a)} - 1\right)
\label{eq:mu_ij_derived}
\end{equation}
It is clear from Eq. (\ref{eq:Bi}) that if $s = a$, $B_i(a, a) = B_j(a, a) = 1$,
and $\mu_{ij}^{(a)} = 0$, as expected.

It is convenient to define
\begin{equation}
\delta f = f_s - f_a
\label{eq:deltaf}
\end{equation}
With this definition, Eq. (\ref{eq:Bi}) becomes
\begin{equation}
B_i(s, a) = 
\frac{g_i(2, s)}{g_i(2, a)}\left(1 - \frac{\delta f}{1 - f_a}\right)
+
A_i(a)\frac{X_i(1)}{X_i(a)}
\left\{
\frac{g_i(1, s)}{g_i(1, a)}\left(1 + \frac{\delta f}{f_a}\right)
- \frac{g_i(2, s)}{g_i(2, a)}\left(1 - \frac{\delta f}{1 - f_a}\right)\right\}
\label{eq:Bi2}
\end{equation}
In the case that there is no species selectivity in forming standard $a$,
$g_i(r, a) = 1$ for all $i$ and $r$.  With this approximation,
$G(1, a) = G(2, a) = 1$, $A_i(a) \to f_a$, and
Eq. (\ref{eq:Bi}) becomes
\begin{equation}
B_i(s, a) = 
g_i(2, s)\left(1 - \frac{\delta f}{1 - f_a}\right)
+
f_a \frac{X_i(1)}{X_i(a)}
\left\{
g_i(1, s)\left(1 + \frac{\delta f}{f_a}\right)
- g_i(2, s)\left(1 - \frac{\delta f}{1 - f_a}\right)\right\}
\label{eq:Bi3}
\end{equation}
If there is also no species selectivity in forming sample $s$, then
$g_i(r, s) = 1$ for all $i$ and $r$ and
Eq. (\ref{eq:Bi}) becomes
\begin{equation}
B_i(s, a) = 1 + \frac{\delta f}{1 - f_a} \left(\frac{X_i(1)}{X_i(a)} - 1\right)
\label{eq:Bi4}
\end{equation}

\section{Simple Model for the NC-CC Dichotomy}

A plausible model for the origin of the NC-CC dichotomy is that CC
region in the early Solar System  contained a higher concentration of NRLEE
dust than the NC region.  This could occur if NRLEE dust preferentially
aggregated into larger dust grains in the NC region and then largely
spiraled into the Sun.  NC compositions would thus arise from mixes
with a lower concentration of NRLEE dust than CC material.

Consider that $CI$ composition reflects the composition of the molecular
cloud (MC) from which the Solar System formed; thus, let $a = CI = MC$ and
$f_{CI}$ is the mass fraction of NRLEE dust in the molecular cloud dust.
Reservoir 1 is the NRLEE dust (denoted $N$) and reservoir 2 is the
complement to the NRLEE dust (denoted $\bar{N}$),
that is, all the dust that is not NRLEE dust.
With the assumption that there is no species selectivity in forming a
sample $s$ with a different concentration $f_s$ of NRLEE dust, one finds
from Eq. (\ref{eq:Bi4})
\begin{equation}
B_i(s, CI) = 1 + \frac{\delta f}{1 - f_{CI}}
\left(\frac{X_i(N)}{X_i(MC)} - 1\right)
\label{eq:BiCI}
\end{equation}
with $\delta f = f_s - f_{CI}$.  A similar equation holds for species $j$.
An anomaly relative to $CI$ as the standard is then,
by Eq. (\ref{eq:mu_ij_derived_ci}),
\begin{equation}
\mu_{ij}^{(CI)}(s) = 10^6 \left(\frac{B_i(s, CI)}{B_j(s, CI)} - 1\right)
\label{eq:mu_ij_derived_ci}
\end{equation}
By Eq. (\ref{eq:diff_mu4}),
the anomaly relative to a terrestrial standard ($b = \oplus$) is 
\begin{equation}
\mu_{ij}^{(\oplus)}(s) = 10^6 \left(\frac{B_i(s, CI)}{B_j(s, CI)} - 1\right)
\left(1 + \frac{\mu_{ij}^{(\oplus)}(CI)}{10^6}\right)
+ \mu_{ij}^{(\oplus)}(CI)
\label{eq:mu_ij_derived_earth}
\end{equation}

A few reasonable approximations yield further insight.  First, the GCE
models in this work indicate $f_{CI} \approx 10^{-4}$; thus, $f_{CI} \ll 1$
and may be neglected in Eq. (\ref{eq:BiCI}).
Second, terms including $\delta f$ are small and may
be expanded in a binomial approximation ($(1 + x)^p \approx 1 + px$ for
$\vert x \vert < 1$).
Finally, for species of interest in this work,
$\mu_{ij}^{(\oplus)}(CI) \approx 200$, so
$\mu_{ij}^{(\oplus)}(CI)/10^6 \approx 2\times 10^{-4}$,
which may be neglected compared to unity.
Thus, Eq. (\ref{eq:mu_ij_derived_earth}) becomes
\begin{equation}
\mu_{ij}^{(\oplus)}(s) \approx \delta f_6
\left(\frac{X_i(N)}{X_i(MC)} - \frac{X_j(N)}{X_j(MC)}\right)
+ \mu_{ij}^{(\oplus)}(CI)
\label{eq:mu_ij_derived_earth2}
\end{equation}
where $\delta f = \delta f_6 \times 10^{-6}$;
that is, $\delta f_6$ is the ppm deviation in
the contribution of NRLEE dust to the sample compared to the NRLEE dust
contribution to CI.  In this model, then, anomalies scale close to linearly
with varying contribution of NRLEE dust.  Furthermore, if NRLEE dust is
more enriched in species $i$ than species $j$ relative to the average
Solar composition ($MC$), then samples depleted in NRLEE dust
($\delta_6 < 0$) will show anomalies closer to NC than CC composition.
For example, for GCE models in this work, $X_i(N)/X_i(MC) \approx 1000$
for $i = {^{50}{\rm Ti}}$ and $X_j(N)/X_j(MC) \approx 20$ for
$j = {^{47}{\rm Ti}}$;
thus, since $\mu^{50}{\rm Ti}(CI) \approx 200$, $\delta f_6 \approx -0.2$
to obtain the Earth's composition relative to $CI$.
Since $f_{CI} \approx 10^{-4}$,
$\delta f_6 = -0.2$ corresponds to an approximate 2 per mil
deficit in NRLEE in NC material relative to $CI$.

\section{CAI-Like Material}

Concentration of refractory material from NRLEE and complementary dust
gives rise to CAI-like material.  In the present model, the mass fraction
of a species $i$ in the CAI-like material relative to CI is,
from Eq. (\ref{eq:Xi}),
\begin{equation}
\frac{X_i(CAI)}{X_i(CI)} = N(CAI, CI) B_i(CAI, CI) 
\label{eq:Xi_cai}
\end{equation}
If the fraction of NRLEE dust going into the CAI-like material is
$f_{CAI}$ and $\delta f = f_{CAI} - f_{CI}$, then Eq. (\ref{eq:Bi3})
becomes
\begin{equation}
B_i(CAI, CI) = \\  
g_i(\bar{N}, CAI)\left(\frac{1 - f_{CAI}}{1 - f_{CI}}\right)
\left(1 + C_i \frac{X_i(N)}{X_i(MC)}\right)
\label{eq:BiCAI}
\end{equation}
where
\begin{equation}
C_i = f_{CAI} \left[\frac{g_i(N, CAI)}{g_i(\bar{N}, CAI)}
\left(\frac{1 - f_{CI}}{1 - f_{CAI}}\right) - \frac{f_{CI}}{f_{CAI}}\right]
\label{eq:C}
\end{equation}

An anomaly in the CAI-like material relative to CI is
\begin{equation}
\mu_{ij}^{(CI)}(CAI) = 10^6 \left(\frac{X_i(N)/X_i(MC)}{X_j(N)/X_j(MC)} - 1
\right) = 10^6 \left(\frac{B_i(CAI, CI) - B_j(CAI, CI)}{B_j(CAI, CI)}\right)
\label{eq:mu_cai}
\end{equation}
Now with the assumption of no isotope selectivity for the same element,
$g_i(r, CAI) = g_j(r, CAI)$ for $r = N$ or $\bar{N}$; hence, $C_i = C_j$.
In this case,
\begin{equation}
\mu_{ij}^{(CI)}(CAI) = 10^6 C_i
\left(\frac{X_i(N)}{X_i(MC)} - \frac{X_j(N)}{X_j(MC)}\right)
\left(1 + C_i \frac{X_j(N)}{X_j(MC)}\right)^{-1}
\label{eq:mu_cai_prop}
\end{equation}
and, if $X_i(N)/X_i(MC) > X_j(N)/X_j(MC)$,
 the sign of the anomaly relative to CI will be given by the sign of $C_i$.
Since $f_{CI} \ll 1$ and $f_{CAI} \ll 1$,
$\mu_{ij}^{(CI)}(CAI) < 0$ for
\begin{equation}
\frac{g_i(N, CAI)}{g_i(\bar{N}, CAI)} < \frac{f_{CI}}{f_{CAI}}
\label{eq:ratio_cai}
\end{equation}

\section{A Simple Model for Compositional Variation in CC Materials}

A model for the compositional variation in CC materials is that CAI-like
matter, formed in the inner Solar disk, gets transported to the CC region
and mixes with the initial molecular cloud material (with CI composition).
Thus, CC material is a mix of reservior 1 (CAI-like matter--denoted CAI)
with CI.
From Eq. (\ref{eq:mu_g2a2}), and with neglect of species selectivity,
the anomaly in CC sample $s$ is
\begin{equation}
\mu_{ij}^{(CI)}(s) =
10^6 \left(\frac{f_s \left[\frac{X_i(CAI)}{X_i(CI)} -
\frac{X_j(CAI)}{X_j(CI)}\right]}
{f_s \frac{X_j(CAI)}{X_j(CI)} + \left(1 - f_s\right)}\right)
\label{eq:mu_g_cai}
\end{equation}
Here $f_s$ is the fraction of the mass of the final CC material contributed
by the CAI-like matter.
This equation may be written
\begin{equation}
\mu_{ij}^{(CI)}(s) =
10^6 \left(\frac{f_{sj}}{1 - f_s + f_{sj}}\right)
\left(\frac{X_i(CAI)/X_i(CI)}{X_j(CAI)/X_j(CI)} - 1\right)
\label{eq:mu_g_cai2}
\end{equation}
where
\begin{equation}
f_{sj} = \frac{X_j(CAI)}{X_j(CI)} f_s
\label{eq:fsj}
\end{equation}
which is now a species specific quantity.
Thus,
\begin{equation}
\mu_{ij}^{(CI)}(s) =
\left(\frac{f_{sj}}{1 - f_s + f_{sj}}\right) \mu_{ij}^{(CI)}(CAI)
\label{eq:mu_g_cai3}
\end{equation}
If $f_s = 0$, there is no contribution
from the CAI-like material, and $\mu_{ij}^{(CI)}(s) = 0$, as expected.
If $f_s = 1$, $\mu_{ij}^{(CI)}(s) = \mu_{ij}^{(CI)}(CAI)$, which makes
sense because the mixture only includes CAI-like matter.
It is also possible for $\mu_{ij}^{(CI)}(s) \approx \mu_{ij}^{(CI)}(CAI)$ even
when $f_s < 1$ if $X_j(CAI) / X_j(CI) \gg 1$ so
that $f_{sj} \gg 1$ because, in this
case, species $j$ is highly enriched in the CAI-like material relative to
CI.

\end{document}

